\clearpage
\SourceLeaf{39}

\begin{center}
\textbf{[MECHANICAL THEOREM ON THE ROOTS OF A POLYNOMIAL
$f(x)$ AND ITS DERIVATIVE $f'(x)$.]}
\end{center}

\noindent\hspace*{\parindent}If $\begin{array}{c}a,\,b,\,c,\,.\,.,\,m,\,n\\ a',\,b',\,c',\,.\,.,\,m'\end{array}$ are the
roots of the equation $\begin{array}{c}f\,x=0\\ f'\,x=0,\end{array}$ where $f'x=\dfrac{d\,f\,x}{d\,x}$,
and if the same letters denote the corresponding points in the plane, then, if
we imagine at $a,\,b,\,c,\,.\,.,\,m,\,n$ equal repelling or attracting masses whose forces vary inversely
with distance, there is equilibrium at $a',\,b',\,c',\,.\,.,\,m'$.

\Needspace{9\baselineskip}
\begin{center}
\textbf{REMARKS ON THE TWO PRECEDING ALGEBRAIC NOTES.}
\end{center}

\noindent\hspace*{\parindent}The first of the two notes given here probably dates from 1826--28. It uses the fifth roots of
unity to solve a system of five linear equations whose determinant has a particularly simple,
symmetric structure. In proving the resulting expressions
for the five quantities $p,\,q,\,r,\,s,\,t$, note
that the norm of the complex number $a+b\varepsilon+c\varepsilon\varepsilon+d\varepsilon^3+e\varepsilon^4$
can be expressed in the form $\frac{1}{4}(5M-m\mathfrak{M})$. 

\textsc{Gauss} probably wrote down the second note in 1846. For the proof,
position the plane so that one of the solutions of $f'(x)=0$ lies at the origin.
Then

\[
\begin{aligned}
\frac{1}{a}+\frac{1}{b}+\frac{1}{c}+\cdots+\frac{1}{n}&=0,\\[3pt]
\frac{a_0}{a\,a_0}+\frac{b_0}{b\,b_0}+\frac{c_0}{c\,c_0}+\cdots+\frac{n_0}{n\,n_0}&=0,
\end{aligned}
\]

\noindent where $a_0$ is the complex conjugate of $a$, $b_0$ of $b$, and so on. The mechanical interpretation of the
last equation yields \textsc{Gauss}’s theorem.
\par\nobreak\hfill\textsc{Fricke.}\par

